\section{Introduction}
As deep neural networks (DNNs) scale and begin to be deployed in increasingly high-stakes settings, it will be imperative to develop a concrete understanding of how these models perform computations and ultimately make decisions.
Towards this end, research in mechanistic interpretability has focused on identifying sub-structures of these models---referred to as \emph{circuits}---and understanding the function and formation of these circuits on the subtasks that they are responsible for.
In order to extract a generalizable understanding of circuits, researchers have formulated a key hypothesis \textbf{universality} \citep{li2015convergent, olah2020zoom}, which suggests that similar networks trained on similar data will form similar circuits. On the other hand, the \textbf{manifold} hypothesis \citep{bengio2013representation, goodfellow2016deep}, suggests that representation learning consists of finding a lower-dimensional manifold for the data. If these hypotheses were to be false, the task of interpreting large scale models becomes dire---there would be little hope of identifying common patterns across initializations, architectures, and datasets.

Yet, recent work claimed totally disparate and disjoint circuits were learned by DNNs trained on \textbf{the exact same data} using modular addition $(a+b)\bmod n = c$ \citep{zhong2023the}. Their results seemingly provide a counter-example to the universality hypothesis, thus suggesting the pursuit of learning interpretable principles that generalize across tasks may be doomed. In fact, this suggests that the identification of simple circuits in larger neural networks could be combinatorially difficult---if different small-scale models trained on one task can learn totally different circuits with no commonality, then large language models (LLMs) may learn many disjoint circuits for a task within their weights simultaneously. 

Our primary contribution to interpreting modular addition, \textit{i.e.} the dataset $(a+b)\bmod n = c$ is the resolution of the \textit{seemingly} disparate circuits found by \cite{zhong2023the}.

\begin{enumerate}
    \item We show that, under certain conditions on the structure of learned embeddings, all networks studied by \citet{zhong2023the} and \citet{mccracken2025uncoveringuniversalabstractalgorithm} learn preactivations having equivalent geometry and topology, which can be expressed in closed form;
    \item These closed-form equations allow us to rigorously claim that all architectures we study universally make use of the same class of manifolds;
    \item We introduce new tools, leveraging topological data analysis, to empirically validate these aforementioned conditions, providing extensive empirical evidence of shared representation geometry across networks.
\end{enumerate}


Altogether, our results restore the possibility that the universality hypothesis is true, as \cite{zhong2023the}'s architectures are no longer shown to be a counter-example. 